\sscp{\tsc{Nuclear Tanimbar-Bomberai}}{11-05-01}
Within \tsc{Tanimbar-Bomberai} there is evidence for a \tsc{Nuclear
Tanimbar-Bomberai} node which contains two further nodes: \tsc{Kei-Fordata}
and \tsc{Bomberai Tip}, the latter of which contains Onin, Sekar, and Uruangnirin.
This structure is different from that proposed by \citet[279]{Blust1993}
who proposed a group which “[{\ldots}] includes Yamdena
and the North Bomberai languages under its lowest node,
joined by Fordata and Kei under the next higher node,”\footnote{
		\citeauthor{Blust1993}'s (\citeyear[273f]{Blust1993})
		North Bomberai group includes Onin,
		Sekar, Uruangnirin, and Arguni.
		Though he notes that the inclusion of Arguni is tentative
		and that it may instead have converged with the other languages
		due to contact.}
Under Blust’s conception Fordata is a primary
node while Yamdena and \tsc{Bomberai Tip} are a separate node.
Instead, we find evidence that Fordata (and Kei) form a single
node with \tsc{Bomberai Tip}
and that Yamdena is a separate node.
There are three sound changes which define \tsc{Nuclear Tanimbar-Bomberai}:
the changes affecting PMP *d,
glide insertion before vowel initial words,
and *j > \it{r} in reflexes of *ŋajan ‘name’,
*huaji ‘younger sibling’, and *suja ‘pitfall trap’ (see \trf{11-22}).

Firstly, the changes affecting PMP *d provide good evidence for
\tsc{Nuclear Tanimbar-Bomberai}, as we can propose PMP *d/*z > **l /{\gap}Vr for \tsc{Nuclear Tanimbar-Bomberai}.
However, before this sound change can be properly understood,
it is necessary to discuss the changes affecting *l in the \tsc{Bomberai Tip} languages.

PMP *l/*n/*d/*z merge in the \tsc{Bomberai Tip} languages.
In most environments the result of this merger is \it{n},
and we can propose *l/*n/*d/*z > **n for Proto-Bomberai Tip.
This is the main defining change of \tsc{Bomberai Tip}.
Sekar and Onin subsequently have **n > \it{r} /{\gap}Vr.
Examples of PMP *l > \it{n} in \tsc{Bomberai Tip} are shown in
\trf{11-25} alongside cognates in other \tsc{Tanimbar-Bomberai} languages for comparison.
Examples of PMP *d/*z > \it{n} in these languages are given in
\trf{11-02} and \trf{11-03} respectively.

\begin{table}\tcp{\tsc{Tanimbar-Bomberai} *l}{11-25}
\begin{adjustbox}{width=\textwidth}
	\st{2pt}\begin{tabular}{llllllll}\lsptoprule
*gloss	&	five	&	laugh	&	body hair	&	skin	&	moon	&	three	&	ear\\
PMP	&	*\tbr{l}ima	&	**ma\tbr{l}ip	&	*bu\tbr{l}u	&	*ku\tbr{l}it	&	*bu\tbr{l}an	&	*tə\tbr{l}u	&	*ta\tbr{l}iŋa\\	\midrule
Yamdena	&	{\hr}lim	&	-malip	&	{\hr}fulu-ŋw	&	{\hr}kulit	&	{\hr}bulan	&	{\hr}tel	&	{\hr}tliŋa-ŋw\\	\midrule
Fordata	&	{\hr}(i)lima	&	{\hr\hr}malit	&	{\hr}vulu-n	&	{\hr}uli-n	&	{\hr}vulən	&	{\hr}telu, itelu	&	\cg{(aru-n)}\\
Kei	&	{\hr}lim	&	{\hr}malit	&	{\hr}vuu-n	&	{\hr}uli-n	&	{\hr}vuan	&	{\hr}tel	&	\cg{(aru-n)}\\
Onin	&	{\hr}\tbr{n}ima	&	{\hr}a-ma\tbr{n}if	&	{\hr}bu\tbr{n}i-n	&	{\hr}ku\tbr{n}it	&	{\hr}pu\tbr{n}an	&	{\hr}te\tbr{n}i	&	{\hr}ta\tbr{n}iga-n\\
Sekar	&	{\hr}\tbr{n}ima	&	{\hr}a-ma\tbr{n}if	&	{\hr}bu\tbr{n}i-m	&	{\hr}u\tbr{n}it	&	{\hr}bu\tbr{n}an	&	{\hr}te\tbr{n}i	&	{\hr}ta\tbr{n}iga-n\\
Uruang.	&	{\hr}\tbr{n}ima	&	{\hr}ma\tbr{n}if	&	{\hr}bu\tbr{n}in	&	{\hr}u\tbr{n}it	&	{\hr}pu\tbr{n}an	&	{\hr}te\tbr{n}i	&	{\hr}ta\tbr{n}iŋan\\	\midrule
Arguni	&	{\hr}rím	&	{\hr}m\<y\>árif	&	{\hr}pu-púre	&	{\hr}árut	&	{\hr}púrin	&	{\hr}táòr	&	{\hr}tigyánáne\\
Bedoanas	&	{\hr}rim	&	{\hr}m\<i\>ári	&	{\hr}urura	&	{\hr}arut	&	{\hr}úrin	&	{\hr}taur	&	\\
Erokw.	&	{\hr}rim	&	{\hr}morif	&	{\hr}pu-pura	&	{\hr}aru	&	{\hr}púrin	&	{\hr}tour	&	{\hr}tiʔinyanya\\
Irarutu	&	\cg{(fradfid)}	&	{\hr}brif	&	{\hr}fru	&	{\hr}rit	&	\cg{(seba)}	&	{\hr}tor	&	{\hr}tgra\\
Kuri	&	\cg{(bradwid)}	&	{\hr}itə-briβ	&	{\hr}ˈβaru, baru	&	{\hr}ˈriːtə̆	&	\cg{(sarabɛ)}	&	{\hr}tɔr	&	{\hr}tə̆gra\\
		\lspbottomrule
	\end{tabular}
\end{adjustbox}
\end{table}

Examples of *n/*l > \it{r} /{\gap}V\it{r} in Onin
and Sekar are given in \trf{11-26}.
Note particularly that this change affects PMP *n,
as seen in *niuR > Onin, Sekar \it{rur} ‘coconut’.
Uruangnirin has *n/*l > \it{n} in the same environment.
Both these facts indicate that *l > \it{r} in Onin
and Sekar occurred after *l > **n in Proto-Bomberai Tip.
Thus, Onin and Sekar share *l > **n > \it{r} /{\gap}V\it{r}.

\begin{table}\tcp{Onin and Sekar **n > \it{r} /{\gap}V\it{r}}{11-26}
%\begin{adjustbox}{width=\textwidth}
	\begin{threeparttable}\st{7pt}
	\begin{tabular}{llll}\lsptoprule
*gloss	&	coconut	&	day, sun	&	voice, language\\
PMP	&	*niuR	&	*qaləjaw	&	**liRi\su{†}\\	\midrule
Yamdena	&	{\hr}nur-e	&	{\hr}ler-e	&	{\hr}liri-ŋw\\	\midrule
Fordata	&	{\hr}nuur	&	{\hr}lera	&	\\
Kei	&	{\hr}nuur	&	{\hr}ler	&	\cg{(ŋrihi)}\\
Onin	&	{\hr}\tbr{r}ur	&	{\hr}\tbr{r}era	&	{\hr}\tbr{r}iri-r\\
Sekar	&	{\hr}\tbr{r}ur	&	{\hr}\tbr{r}era	&	{\hr}\tbr{r}iri-r\\
Uruang.	&	\cg{(wat)}	&	{\hr}\tbr{n}era	&	{\hr}\tbr{n}iri-n\\	\midrule
Arguni	&	{\hr}náwur	&	\cg{(rúas)}	&	\cg{(ariga-m)}\\
Bedoanas	&	{\hr}nau	&	\cg{(sáni)}	&	\\
Erokw.	&	{\hr}nau	&	\cg{(úruas)}	&	\cg{(ariʔan)}\\
Irarutu	&	\cg{(umagi)}	&	{\hr}re, rre	&	{\hr}rü\\
Kuri	&	\cg{(ˈakadi)}	&	{\hr}ror, rərɔːr	&	\cg{(garnɛ, snan)}\\
		\lspbottomrule
	\end{tabular}
		\begin{tablenotes}\tnsize
			\item[†] {The Kei,
								Arguni and Erokwanas words for ‘voice’
								cannot be regularly derived from putative PCMP *liRi.
								Medial \it{h} in Kei points to earlier **s.
								Arguni \it{g} Erokwanas \it{ʔ} point to earlier **ŋ.}
		\end{tablenotes}
	\end{threeparttable}
%\end{adjustbox}
\end{table}

With these sound changes in mind,
we now turn to the changes affecting PMP *d/*z.
Kei and Fordata have *d > \it{l} /{\gap}Vr
and *d > \it{r} elsewhere.
PMP *d merges with *l/*z/*n in \tsc{Bomberai Tip} -- that is
*d/*l/*z > **n > \it{r} /{\gap}Vr in Onin
and Sekar, and *d/*l/*z > \it{n} elsewhere in all three languages.
Examples are shown in \trf{11-27}.
(*diRi ‘stand’ retains *d = \it{d} in \tsc{Kei-Fordata},
perhaps due to an earlier stative prefix *ma- which is preserved
in Yamdena and \tsc{Bomberai Tip}.)

To account for these changes,
we can posit that *d/*z > **l /{\gap}V\it{r} had occurred by
the time of Proto-Nuclear Tanimbar-Bomberai,
where conditioning \it{r} is from *R/*j/*d/*z.
After the split between \tsc{Kei-Fordata}
and \tsc{Bomberai Tip}, this **l then merged with *l/*d/*z/*n in \tsc{Bomberai Tip}.
Thus, shared *d > **l /{\gap}V\it{r} provides evidence for \tsc{Nuclear Tanimbar-Bomberai}.

Supporting evidence for \tsc{Nuclear Tanimbar-Bomberai} comes
from insertion of \it{y} before certain vowel-initial words,
particularly in vowels beginning with \it{{\#}a} in other languages (\srf{13-02-02}).\footnote{
		Many languages in several subgroups in the region have \it{y}-insertion
		before nouns beginning with \it{{\#}a} (along with other irregular
		onsets; see data and discussion in \srf{13-02-02})
		and thus this is not an exclusively shared innovation.
		In many other areas the phenomenon is sporadic
		and does not have subgrouping value,
		but here the consistency of the pattern means that it can be
		taken as supporting evidence for \tsc{Nuclear Tanimbar-Bomberai}.}
Examples are given in \trf{11-28}.
This change also occurs in some Yamdena
and \tsc{East Bomberai} words,
but not to the same extent as in \tsc{Nuclear Tanimbar-Bomberai}.
Uruangnirin has undergone subsequent *y > \it{l},
also seen in *buqa\tbr{y}a > \it{pua\tbr{l}a} ‘crocodile’,
and *i-aku > **\tbr{y}aku > \it{\tbr{l}au}.

\begin{table}\tcp{\tsc{Tanimbar-Bomberai} *d}{11-27}
\begin{adjustbox}{width=\textwidth}
	\begin{threeparttable}\st{2pt}
	\begin{tabular}{lllll|llll}\lsptoprule
*gloss	&	two	&	inside	&	leaf	&	behind	&	thorn\su{†}	&	blood	&	bathe	&	stand\\
PMP	&	*\tbr{d}uha		&	*\tbr{d}aləm	&	*\tbr{d}ahun	&	*ma-u\tbr{d}əhi	&	*\tbr{d}uRi	&	*\tbr{d}aRaq	&	*\tbr{d}iRus	&	*ma-\tbr{d}iRi\su{‡}\\	\midrule
Yamdena	&	{\hr}du	&	{\hr}dalam	&	{\hr}don-e	&	{\hr}mudy	&	{\hr}duri	&	{\hr}dar-e	&	{\hr}diris	&	{\hr}namdiry\\
Fordata	&	{\hr}\tbr{r}ua	&	{\hr}\tbr{r}ala-n	&	{\hr}\tbr{r}oa	&	{\hr}mu\tbr{r}i-n	&	{\hr}\tbr{l}uri-n	&	{\hr}\tbr{l}ara-n	&	{\hr}\tbr{l}iru	&	{\hr}n-\tbr{d}iri\\
Kei	&	{\hr}\tbr{r}u	&	{\hr}\tbr{r}aa-n	&	{\hr}\tbr{r}on	&	{\hr}mu\tbr{r}	&	{\hr}\tbr{l}uri-n	&	{\hr}\tbr{l}ar	&	{\hr}n-\tbr{l}iruk	&	{\hr}\tbr{d}ir\\
Onin	&	{\hr}\tbr{n}ua	&	{\hr}\tbr{n}anam	&	{\hr}wi \tbr{n}on-in	&	{\hr}mu\tbr{n}i	&	{\hr}\tbr{r}uri-r	&	{\hr}\tbr{r}ara	&	\cg{(orat)}	&	{\hr}ma\tbr{r}iri\\
Sekar	&	{\hr}\tbr{n}ua	&	{\hr}\tbr{n}anam	&	\cg{(kafak-in)}	&	\cg{(tabi-n)}	&	{\hr}\tbr{r}uri-r	&	{\hr}\tbr{r}ara	&	\cg{(a-irit)}	&	{\hr}m\tbr{r}iri\\
Uruang.	&	{\hr}\tbr{n}ua	&	{\hr}\tbr{n}anam	&	{\hr}\tbr{n}on-in	&	{\hr}mu\tbr{n}in	&	{\hr}\tbr{n}uri-n	&	{\hr}\tbr{n}ara	&	\cg{(li-sofa)}	&	{\hr}a-m\tbr{r}iri\\	\midrule
Arguni	&	{\hr}ru	&	{\hr}ramo-n	&	\cg{(náne)}	&		&	\cg{(tór)}	&	{\hr}ráre	&	{\hr}a-resə	&	\cg{(gáser)}\\
Bedoanas	&	{\hr}ruyu	&		&	{\hr}rína	&		&	\cg{(tóʔor)}	&	\cg{(rita)}	&		&	\cg{(usar)}\\
Erokw.	&	{\hr}ru	&	{\hr}-ramu	&	\cg{(manda-nánia)}	&	\cg{(toʔar)}	&	\cg{(rita)}	&	{\hr}rasyaf	&	\cg{(waser)}\\
Irarutu	&	{\hr}rifo	&	\cg{(gan)}	&	{\hr}ro	&	\cg{(barie)}	&	{\hr}rur	&	\cg{(wams)}	&		&	{\hr}n-marir\\
Kuri	&	{\hr}ru	&	\cg{(gurˈnɛ)}	&	{\hr}ruɛˈrɔ	&		&	{\hr}rur	&	\cg{(wams)}	&		&	{\hr}itə-brir\\
		\lspbottomrule
	\end{tabular}
		\begin{tablenotes}\tnsize
			\item[†]	{PMP *duRi is reconstructed with the meaning ‘thorn, splinter, fish bone’.
								The Yamdena, Sekar, and Onin reflexes mean ‘thorn,
								bone’, while other reflexes are only known to mean ‘bone’.}
			\item[‡]	{The Onin,
								Sekar, Uruangnirin, Arguni, Bedoanas,
								and Erokwanas words were all given in response to ‘get up,
								wake up’ (Indonesian \it{baŋun}).
								The Yamdena, Fordata, and Kei words all mean ‘stand’.}
		\end{tablenotes}
	\end{threeparttable}
\end{adjustbox}
\end{table}

\begin{table}[p]\centering\tcp{\tsc{Tanimbar-Bomberai} \it{y}-insertion}{11-28}
%\begin{adjustbox}{width=\textwidth}
	\begin{threeparttable}\st{2.5pt}
	\begin{tabular}{llllllll}\lsptoprule
*gloss	&	child	&	father	&	liver	&	lime	&	dog	&	fire	&	forest\\
PMP	&	*anak	&	*ama	&	*qatay	&	*qapuR\su{‡}	&	*asu	&	*hapuy	&	*halas\\	\midrule
Yamdena	&	{\hr}  anaky	&	{\hr}  ama-n	&	\cg{(bati)}	&	{\hr}\tbr{y}afur	&	{\hr}  asu	&	{\hr}  au	&	{\hr}  alas\\	\midrule
Fordata	&	{\hr}\tbr{y}ana-n	&	{\hr}\tbr{y}ama-n	&	{\hr}\tbr{y}ata-n	&	{\hr}\tbr{y}afur	&	{\hr}\tbr{y}aha	&	{\hr}\tbr{y}afu	&	{\hr}\tbr{y}alat\su{Δ}\\
Kei	&	{\hr}\tbr{y}ana-n	&	{\hr}\tbr{y}ama-n	&	{\hr}\tbr{y}ata-n\su{†}	&	{\hr}\tbr{y}afur	&	{\hr}\tbr{y}ahaw	&	{\hr}\tbr{y}af	&	{\hr}\tbr{y}aat\\
Onin	&	{\hr}\tbr{y}ana-n	&	{\hr}\tbr{y}ama-n	&	{\hr}\tbr{y}ata-n	&	\cg{(lofin)}	&	\cg{(agia)}	&	{\hr}\tbr{y}afi	&	\cg{(sananan)}\\
Sekar	&	\cg{(kukanak)}	&	{\hr}\tbr{y}ama-n	&	\cg{(nawan)}	&	{\hr}\tbr{y}afer	&	{\hr}\tbr{y}asi	&	{\hr}\tbr{y}afi	&	\cg{(kenin)}\\
Uruang.	&	\cg{(keka)}	&	{\hr}\tbr{l}ama-n	&	{\hr}\tbr{l}ata-n	&	{\hr}\tbr{l}afur	&	{\hr}\tbr{l}asi	&	{\hr}\tbr{l}afi	&	\cg{(unanan)}\\	\midrule
Arguni	&	\cg{(kákon)}	&	{\hr}támana	&		&	{\hr}afir	&	\cg{(ófun)}	&	{\hr}\tbr{iy}áf	&	\cg{(áter)}\\
Bedoanas	&	\cg{(kókon)}	&	\cg{(máma)}	&	\cg{(tapóti)}	&		&	\cg{(áfun)}	&	{\hr}\tbr{í}af	&	\\
Erokw.	&	\cg{(kakon)}	&	\cg{(bápa)}	&	\cg{(wade)}	&	{\hr}afir	&	\cg{(áun)}	&	{\hr}\tbr{í}af	&	\cg{(waríri)}\\
%Irarutu	&	\cg{(mo)}	&	\cg{(anʤe)}	&	{\hr}ti	&		&	\cg{(fun)}	&	\cg{(fena)}	&	\cg{(witu)}\\
%Kuri	&	\cg{(yenatu)}	&		&		&	\cg{(siˈru)}	&	\cg{(abun)}	&	\cg{(wɛta)}	&	\cg{(utˈrova)}\\
		\lspbottomrule
	\end{tabular}
		\begin{tablenotes}\tnsize
			\item[†]	{Kei \it{yata-n} is given meaning ‘heart’.}
			\item[‡]	{Onin \it{lofin} ‘mineral lime’ may ultimately be from *qapuR
								via an intermediate language.
								Uruangnirin \it{lafur} ‘mineral lime’ has irregular final *u
								= \it{u} rather than regular *u > \it{i} /σ{\gap}.}
			\item[Δ]	{Fordata \it{yalat} means ‘dense forest,
								stand of trees’.}
		\end{tablenotes}
	\end{threeparttable}
%\end{adjustbox}
\end{table}

Within \tsc{Nuclear Tanimbar-Bomberai},
Kei and Fordata can be grouped together on the basis of the following
innovations: *b > \it{v},
*ŋ > \it{n}, and *s > **h (with subsequent **h > {\0} in some Fordata words).
Two examples of each of these sound changes are given in \trf{11-29},
along with their cognates in other \tsc{Tanimbar-Bomberai} languages
for comparison.

\begin{table}[p]\centering\tcp{\tsc{Kei-Fordata} *b > \it{v}; *s > \it{h}; *ŋ > \it{n}}{11-29}
%\begin{adjustbox}{width=\textwidth}
	\begin{threeparttable}\st{4pt}
	\begin{tabular}{lll|ll|ll}\lsptoprule
*gloss	&	stone						&	new											&	seawater				&	flesh							&	sky								&	weep\\
PMP			&	*\tbr{b}atu			&	*\tbr{b}aqəRu						&	*ta\tbr{s}ik		&	*i\tbr{s}i				&	*la\tbr{ŋ}it			&	*ta\tbr{ŋ}is\su{‡}		\\	\midrule
Yamdena	&	{\hr}\tbb{b}aty	&	{\hr}be{\tl}\tbb{b}ery	&	{\hr}ta\tbb{s}ik&	{\hr}i\tbb{s}in		&	{\hr}la\tbb{ŋ}it	&	{\hr}-tasi\tbb{ŋ}			\\	\midrule
Fordata	&	{\hr}\tbr{v}atu	&	{\hr}\tbr{v}eru					&	{\hr}ta\tbr{h}at&	{\hr}i\tbr{h}i-n	&	{\hr}la\tbr{n}it	&	{\hr}n-ta\tbr{n}it		\\
Kei			&	{\hr}\tbr{v}at	&	{\hr}\tbr{v}irun\su{†}	&	{\hr}ta\tbr{h}it&	{\hr}i\tbr{h}i-n	&	{\hr}la\tbr{n}it	&	{\hr}tna\tbr{n}it			\\
Onin		&	{\hr}\tbb{p}ati	&	{\hr}per-\tbb{p}eri			&	{\hr}ta\tbb{s}i	&	{\hr}is{\tl}isin	&	\cg{(laŋit)}			&	\cg{(a-gaga)}					\\
Sekar		&	{\hr}\tbb{b}ati	&	{\hr}\tbb{b}eri					&	\cg{(amirór)}		&	{\hr}is{\tl}isin	&	\cg{(laŋit)}			&	{\hr}a-ta\tbb{g}is		\\
Uruang.	&	{\hr}\tbb{p}atin&	{\hr}per{\tl}\tbb{p}eri	&	{\hr}ta\tbb{s}ik&	{\hr}i\tbb{s}in		&	{\hr}na\tbb{g}it	&	\cg{(a-rekar)}				\\	\midrule
Arguni	&	{\hr}\tbb{p}uát	&	{\hr}pu{\tl}\tbb{p}uaru	&	\cg{(fayóret)}	&	\cg{(sumbuáne)}		&	{\hr}rá\tbb{g}it	&	{\hr}t\<y\>á\tbb{g}īs	\\
Bedoanas&	{\hr}uat				&	{\hr}ur{\tl}uri					&	{\hr}tá\tbb{s}i	&	\cg{(sunuane)}		&	{\hr}ra\tbb{ʔ}it	&	{\hr}t\<i\>á\tbb{ʔ}is	\\
Erokw.	&	{\hr}\tbb{p}uat	&	{\hr}pu{\tl}\tbb{p}uari	&	\cg{(fayórat)}	&	\cg{(sobuana)}		&	{\hr}rá\tbb{ʔ}it	&	{\hr}t\<i\>á\tbb{ʔ}is	\\
Irarutu	&	\cg{(kami)}			&	\cg{(bunat)}						&	\cg{(werfun)}		&	\cg{(ʤe)}					&	{\hr}ra\tbb{g}t		&	{\hr}n-ta\tbb{g}			\\
%Kuri		&	\cg{(kɔbaːb)}		&	\cg{(aˈbɔb)}						&	\cg{(ruarə̆nə̆)}	&	\cg{(yɛ)}					&										&	{\hr}itə-ta\tbb{g}		\\
		\lspbottomrule
	\end{tabular}
		\begin{tablenotes}\tnsize
			\item[†]	{Kei \it{virun} is glossed ‘young (person)’.}
			\item[‡]	{Kei \it{tnanit} means ‘song of farewell for journey or bride’s
								departure’, Fordata \it{n-tanit} is ‘wail, mourn’.
								Yamdena \it{-tasiŋ} has metathesis of the medial and final consonants.}
		\end{tablenotes}
	\end{threeparttable}
%\end{adjustbox}
\end{table}

Alongside \tsc{Kei-Fordata},
the other node of \tsc{Nuclear Tanimbar-Bomberai} is \tsc{Bomberai Tip},
defined by merger of *l/*z/*d/*n
> **n, as discussed above.
In addition, these languages usually share *ŋ > \it{g},
\it{ŋg} (see \trf{11-31} below)
and final *u > \it{i}.
Neither are exclusively shared innovations,
as *ŋ > \it{g} also occurs in \tsc{East Bomberai} (\srf{11-05-02})
and Kowiai (\srf{11-03-02}), while *u > \it{i} also occurs in
Yamdena.
